\section{Theorem-like Environments}

\begin{ruletable}
Definition & Precise meaning of a term, object, property, or construction. \\
Notation & Introduction or declaration of notation. \\
Theorem & Principal mathematical result. \\
Proposition & Useful result of more limited scope. \\
Lemma & Supporting result, usually for a later proof. \\
Corollary & Result following readily from a preceding result. \\
Conjecture & Statement proposed to be true but not proved in full generality. \\
Hypothesis & Named assumption or proposed statement. \\
Formula & Distinguished computational identity or closed form. \\
Algorithm & Explicit step-by-step procedure. \\
Fact & Verified statement not requiring theorem status. \\
Example & Illustration of a definition, construction, phenomenon, or result. \\
Counterexample & Example disproving a general assertion or showing a hypothesis is necessary. \\
Remark & Supplementary context, warning, or qualification. \\
Problem & Mathematical task intended to be solved. \\
Open Problem & Problem not known to be solved in the stated form. \\
Idea & Decisive conceptual insight. \\
Strategy & Broad problem-solving plan. \\
Tactic & Local step within a larger strategy. \\
\end{ruletable}

Proof labels are written as \textit{Proof} and solution labels as \textit{Solution}, each followed by a period.

\begin{notationtable}
$\square$ & Ending for ordinary proofs and solutions. \\
$\blacksquare$ & Ending for substantial, difficult, or particularly important proofs. \\
\textsc{Q.E.D.} & Ending for exceptionally difficult, significant, or paper-level proofs. \\
\end{notationtable}

\Needspace{9\baselineskip}
\noindent\textbf{Enumeration and Subdivision.}
Enumeration labels are set in upright Roman type. Their form is determined by structural role rather than nesting depth.

\begin{ruletable}
\textup{1.1, 1.2, 1.3, \ldots} & Independent child problems or standalone units under a common parent; each is a distinct subproblem that may be stated, solved, and referenced separately. \\
\textup{(1), (2), (3), \ldots} & Ordinary parallel enumeration, such as choices, cases, possibilities, alternatives, examples, or procedural steps; not a subdivision of a single problem or a clause of a formal mathematical statement. \\
\textup{(a), (b), (c), \ldots} & Connected subparts of a single problem or task that share a common statement, setup, or context. \\
\textup{(i), (ii), (iii), \ldots} & Formal mathematical clauses within definitions, theorems, propositions, lemmas, and related statements, such as hypotheses, conditions, properties, conclusions, or equivalent statements. \\
\end{ruletable}

The distinction between \textup{1.1} and \textup{(a)} is independence: \textup{1.1} denotes a distinct child problem under a common parent, whereas \textup{(a)} denotes a connected part of one problem sharing its setup. The distinction between \textup{(1)} and \textup{(i)} is function: \textup{(1)} is used for ordinary parallel enumeration, whereas \textup{(i)} is reserved for clauses within formal mathematical statements. Nested enumeration preserves these structural roles rather than following a fixed hierarchy of label styles.

\Needspace{10\baselineskip}
\noindent\textbf{Mathematical Shorthand.}
Typeset abbreviations may be used in typeset mathematical writing as well as handwritten work. Handwritten shorthand is reserved for notes, scratch work, and blackboard writing.

\Needspace{6\baselineskip}
\noindent\textit{Typeset Abbreviations.}
\begin{notationtable}
\textnormal{i.e.} & ``that is.'' \\
\textnormal{i.o.w.} & ``in other words.'' \\
\textnormal{s.t.} & ``such that''; preferred when an abbreviation is natural in typeset mathematical writing. \\
\textnormal{iff} & ``if and only if.'' \\
\textnormal{w.r.t.} & ``with respect to.'' \\
\textnormal{const} & ``constant.'' \\
\addlinespace[0.35em]
\textnormal{N.t.} & ``note that.'' \\
\textnormal{w.l.o.g.} & ``without loss of generality.'' \\
\textnormal{w.t.s.} & ``want to show (that).'' \\
\textnormal{e.t.s.} & ``enough to show (that).'' \\
\textnormal{t.f.a.e.} & ``the following are equivalent.'' \\
\textnormal{i.t.s.w.} & ``in the same way.'' \\
\textnormal{o.t.o.h.} & ``on the other hand.'' \\
\addlinespace[0.35em]
\textnormal{IVT} & Intermediate Value Theorem. \\
\textnormal{EVT} & Extreme Value Theorem. \\
\textnormal{MVT} & Mean Value Theorem. \\
\textnormal{FTC} & Fundamental Theorem of Calculus. \\
\addlinespace[0.35em]
\textnormal{ODE} & Ordinary Differential Equation. \\
\textnormal{PDE} & Partial Differential Equation. \\
\textnormal{IVP} & Initial Value Problem. \\
\textnormal{BVP} & Boundary Value Problem. \\
\end{notationtable}

\Needspace{6\baselineskip}
\noindent\textit{Handwritten Shorthand.}
\begin{notationtable}
\textnormal{Def} & ``definition.'' \\
\textnormal{Thm} & ``theorem.'' \\
\textnormal{Prop} & ``proposition.'' \\
\textnormal{Lem} & ``lemma.'' \\
\textnormal{Cor} & ``corollary.'' \\
\textnormal{Pf} & ``proof.'' \\
\textnormal{Sol} & ``solution.'' \\
\textnormal{Ex} & ``example.'' \\
\textnormal{Exer} & ``exercise.'' \\
\textnormal{Rmk} & ``remark.'' \\
\textnormal{Sps} & ``suppose (that).'' \\
\addlinespace[0.35em]
\textnormal{cont} & ``continuous.'' \\
\textnormal{cont.ly} & ``continuously.'' \\
\textnormal{difb} & ``differentiable.'' \\
\textnormal{difb.ly} & ``differentiably.'' \\
\textnormal{difl} & ``differential.'' \\
\textnormal{bdd} & ``bounded.'' \\
\textnormal{unif} & ``uniform.'' \\
\textnormal{unif.ly} & ``uniformly.'' \\
\textnormal{lin} & ``linear.'' \\
\textnormal{lin.ly} & ``linearly.'' \\
\textnormal{dep} & ``dependent.'' \\
\textnormal{indep} & ``independent.'' \\
\textnormal{approx} & ``approximate.'' \\
\textnormal{triv} & ``trivial.'' \\
\addlinespace[0.35em]
\textnormal{fnc} & ``function.'' \\
\textnormal{seq} & ``sequence.'' \\
\textnormal{eqn} & ``equation.'' \\
\textnormal{eqlt} & ``equality.'' \\
\textnormal{ineqlt} & ``inequality.'' \\
\textnormal{sys} & ``system.'' \\
\textnormal{comb} & ``combination.'' \\
\textnormal{vec} & ``vector.'' \\
\textnormal{mat} & ``matrix.'' \\
\textnormal{col} & ``column.'' \\
\textnormal{pt} & ``point.'' \\
\textnormal{bdy} & ``boundary.'' \\
\textnormal{nbhd} & ``neighborhood.'' \\
\textnormal{itvl} & ``interval.'' \\
\textnormal{op} & ``operation.'' \\
\textnormal{num} & ``number.'' \\
\textnormal{val} & ``value.'' \\
\textnormal{elmt} & ``element.'' \\
\addlinespace[0.35em]
\textnormal{cw} & ``clockwise.'' \\*
\textnormal{ccw} & ``counterclockwise.'' \\*
\textnormal{w/} & ``with.'' \\*
\textnormal{w/o} & ``without.'' \\
\end{notationtable}

\Needspace{7\baselineskip}
\noindent\textbf{Mathematical Phrases.}
\begin{phrasetable}
\textnormal{We find by inspection that \ldots} & Used when a conclusion follows directly from inspection. \\
\end{phrasetable}

